% (c) Nikita Lisitsa, lisyarus@gmail.com, 2026

\documentclass[10pt]{beamer}

\usepackage[T2A]{fontenc}
\usepackage[russian]{babel}
\usepackage{minted}

\usepackage[most]{tcolorbox}
\tcbuselibrary{minted}

\usepackage{graphicx}
\graphicspath{ {./images/} }

\usepackage{adjustbox}

\usepackage{color}
\usepackage{soul}
\usepackage{multicol}
\usepackage{hyperref}
\usepackage{tikz}

\newcommand{\shrug}[1][]{%
\begin{tikzpicture}[baseline,x=0.8\ht\strutbox,y=0.8\ht\strutbox,line width=0.125ex,#1]
\def\arm{(-2.5,0.95) to (-2,0.95) (-1.9,1) to (-1.5,0) (-1.35,0) to (-0.8,0)};
\draw \arm;
\draw[xscale=-1] \arm;
\def\headpart{(0.6,0) arc[start angle=-40, end angle=40,x radius=0.6,y radius=0.8]};
\draw \headpart;
\draw[xscale=-1] \headpart;
\def\eye{(-0.075,0.15) .. controls (0.02,0) .. (0.075,-0.15)};
\draw[shift={(-0.3,0.8)}] \eye;
\draw[shift={(0,0.85)}] \eye;
% draw mouth
\draw (-0.1,0.2) to [out=15,in=-100] (0.4,0.95); 
\end{tikzpicture}}

\usetheme{metropolis}

\definecolor{red}{rgb}{1,0,0}
\definecolor{codebg}{RGB}{29,35,49}
\setminted{fontsize=\footnotesize}

\makeatletter
\newcommand{\slideimage}[1]{
  \begin{figure}
    \begin{adjustbox}{width=\textwidth, totalheight=\textheight-2\baselineskip-2\baselineskip,keepaspectratio}
      \includegraphics{#1}
    \end{adjustbox}
  \end{figure}
}
\makeatother

\title{Язык программирования C++}
\subtitle{Лекция 14: шаблоны, исключения, раскрутка стека, гарантии при работе с исключениями}
\date{2026}
\author{Никита Лисица}

\setbeamertemplate{footline}[frame number]

\begin{document}

\frame{\titlepage}

\begin{frame}[fragile]
\frametitle{Шаблонные классы}
\begin{itemize}
\usemintedstyle{lightbulb}
\begin{tcblisting}{listing engine=minted, minted language=cpp, minted options={fontsize=\scriptsize}, listing only, colback=codebg, colframe=codebg, rounded corners}
// array.h

template <typename T>
class Array {
  T* data;
  size_t size;

public:
  Array(size_t size)
    : data(new T[size])
    , size(size)
  {}

  ~Array() {
    delete[] data;
  }

  T& operator[] (size_t index) {
    return data[index];
  }
};
\end{tcblisting}
\end{itemize}
\end{frame}

\begin{frame}[fragile]
\frametitle{Шаблонные классы}
\begin{itemize}
\item Реализации методов можно вынести наружу:
\usemintedstyle{lightbulb}
\begin{tcblisting}{listing engine=minted, minted language=cpp, minted options={fontsize=\scriptsize}, listing only, colback=codebg, colframe=codebg, rounded corners}
// array.h

template <typename T>
class Array {
  T* data;
  size_t size;

public:
  Array(size_t size);
  ~Array();
  T& operator[] (size_t index);
};

...
\end{tcblisting}
\end{itemize}
\end{frame}

\begin{frame}[fragile]
\frametitle{Шаблонные классы}
\begin{itemize}
\item Реализации методов можно вынести наружу:
\usemintedstyle{lightbulb}
\begin{tcblisting}{listing engine=minted, minted language=cpp, minted options={fontsize=\scriptsize}, listing only, colback=codebg, colframe=codebg, rounded corners}
// array.h

...

template <typename T>
Array<T>::Array(size_t size)
  : data(new T[size])
  , size(size)
{}

template <typename T>
Array<T>::~Array() {
  delete[] data;
}

template <typename T>
T& Array<T>::operator[] (size_t index) {
  return data[index];
}
\end{tcblisting}
\end{itemize}
\end{frame}

\begin{frame}[fragile]
\frametitle{Шаблонные классы}
\begin{itemize}
\usemintedstyle{lightbulb}
\begin{tcblisting}{listing engine=minted, minted language=cpp, minted options={fontsize=\scriptsize}, listing only, colback=codebg, colframe=codebg, rounded corners}
#include <array.h>

int main() {
  Array<int> array1;
  Array<double> array2;
  Array<Array<int>> array3;
}
\end{tcblisting}
\pause
\usemintedstyle{tango}
\item Подстановку делает \textbf{компилятор}, а не \textbf{препроцессор}!
\pause
\item Код шаблонного класса должен быть в заголовочном файле. Почему? \pause Мы не знаем заранее, какие именно типы могут подставить в качестве \mintinline{cpp}|T|
\pause
\item Методы шаблонного класса всегда \mintinline{cpp}|inline|. Почему? \pause Один и тот же тип \mintinline{cpp}|T| могут подставить в разных cpp-файлах
\end{itemize}
\end{frame}

\begin{frame}[fragile]
\frametitle{Шаблонные классы}
\begin{itemize}
\usemintedstyle{tango}
\item \mintinline{cpp}|T| называется \textit{шаблонным параметром}
\pause
\item \mintinline{cpp}|template <typename T>| и \mintinline{cpp}|template <class T>| ничем не отличаются (лучше \mintinline{cpp}|typename| -- шаблонный параметр не обязан быть классом)
\pause
\item \mintinline{cpp}|template <struct T>| нельзя :)
\pause
\item Класс с конкретными подставленными значениями (например, \mintinline{cpp}|Array<int>|) называется \textit{специализацией (specialization)} шаблона
\pause
\item Можно делать синонимы типов на специализации: \mintinline{cpp}|using IntArray2D = Array<Array<int>>;|
\pause
\item Процесс генерации кода для конкретной специализации называется \textit{инстанциированием (instantiation)} шаблона
\end{itemize}
\end{frame}

\begin{frame}[fragile]
\frametitle{Несколько шаблонных параметров}
\begin{itemize}
\usemintedstyle{lightbulb}
\begin{tcblisting}{listing engine=minted, minted language=cpp, minted options={fontsize=\scriptsize}, listing only, colback=codebg, colframe=codebg, rounded corners}
template <typename Key, typename Value>
class RedBlackTree {
  struct Node {
    Key key;
    Value value;
  };

  Node* root;

public:
  ...
};
\end{tcblisting}
\end{itemize}
\end{frame}

\begin{frame}[fragile]
\frametitle{Значения по умолчанию}
\begin{itemize}
\usemintedstyle{lightbulb}
\begin{tcblisting}{listing engine=minted, minted language=cpp, minted options={fontsize=\scriptsize}, listing only, colback=codebg, colframe=codebg, rounded corners}
template <typename T = float>
struct Vector3 {
  T x, y, z;
};

Vector3 v1; // Vector3<float>
Vector3<int> v2;
Vector3<double> v3;
\end{tcblisting}
\end{itemize}
\end{frame}

\begin{frame}[fragile]
\frametitle{Шаблонные функции}
\begin{itemize}
\usemintedstyle{lightbulb}
\begin{tcblisting}{listing engine=minted, minted language=cpp, minted options={fontsize=\scriptsize}, listing only, colback=codebg, colframe=codebg, rounded corners}
template <typename T>
void swap(T& a, T& b) {
  T temp = a;
  a = b;
  b = temp;
}

int main() {
  int i = 10;
  int j = 20;
  swap<int>(i, j);
}
\end{tcblisting}
\pause
\item Шаблонные функции тоже должны быть в заголовочных файлах!
\end{itemize}
\end{frame}

\begin{frame}[fragile]
\frametitle{Шаблонные функции}
\begin{itemize}
\item Типы аргументов могут зависеть от шаблонных параметров:
\usemintedstyle{lightbulb}
\begin{tcblisting}{listing engine=minted, minted language=cpp, minted options={fontsize=\scriptsize}, listing only, colback=codebg, colframe=codebg, rounded corners}
template <typename T>
void reverse(Array<T>& array)
  size_t size = array.size();
  for (size_t i = 0; i < size / 2; ++i) {
    swap<T>(array[i], array[size - i - 1]);
  }
}

int main() {
  Array<int> array(10);
  ...
  reverse<int>(array);
}
\end{tcblisting}
\end{itemize}
\end{frame}

\begin{frame}[fragile]
\frametitle{Несколько шаблонных параметров}
\begin{itemize}
\usemintedstyle{lightbulb}
\begin{tcblisting}{listing engine=minted, minted language=cpp, minted options={fontsize=\scriptsize}, listing only, colback=codebg, colframe=codebg, rounded corners}
template <typename T, typename H>
void copy(const Array<T>& src, Array<H>& dst)
  dst.resize(src.size());
  for (size_t i = 0; i < src.size(); ++i)
    dst[i] = src[i];
}

int main() {
  Array<int> array(10);
  ...
  Array<double> array2;
  copy<int, double>(array, array2);
}
\end{tcblisting}
\end{itemize}
\end{frame}

\begin{frame}[fragile]
\frametitle{Вывод шаблонных параметров}
\begin{itemize}
\item Часто компилятор может сам вывести шаблонные параметры из типов аргументов:
\usemintedstyle{lightbulb}
\begin{tcblisting}{listing engine=minted, minted language=cpp, minted options={fontsize=\scriptsize}, listing only, colback=codebg, colframe=codebg, rounded corners}
Array<int> array(10);
// Array<T> = Array<int>
// => T = int
reverse(array);

Array<double> array2;
// Array<T> = Array<double>
// => T = int
reverse(double);

// Array<T> = Array<int>
// Array<H> = Array<double>
// => T = int
// => H = double
copy(array, array2);
\end{tcblisting}
\end{itemize}
\end{frame}

\begin{frame}[fragile]
\frametitle{C vs ООП vs Шаблоны}
\begin{itemize}
\usemintedstyle{lightbulb}
\begin{tcblisting}{listing engine=minted, minted language=cpp, minted options={fontsize=\scriptsize}, listing only, colback=codebg, colframe=codebg, rounded corners}
void sort(void* array, size_t size, size_t count, cmp_fn cmp);
\end{tcblisting}
\usemintedstyle{tango}
\item \alert{\textbf{Работает}} для примитивных типов, \textbf{нет} проверок типов, \textbf{нельзя} заинлайнить сравнение
\pause
\usemintedstyle{lightbulb}
\begin{tcblisting}{listing engine=minted, minted language=cpp, minted options={fontsize=\scriptsize}, listing only, colback=codebg, colframe=codebg, rounded corners}
void sort(Comparable** array, size_t size);
\end{tcblisting}
\usemintedstyle{tango}
\item \textbf{Не работает} для примитивных типов, \alert{\textbf{частичная}} проверка типов, \textbf{нельзя} заинлайнить сравнение
\pause
\usemintedstyle{lightbulb}
\begin{tcblisting}{listing engine=minted, minted language=cpp, minted options={fontsize=\scriptsize}, listing only, colback=codebg, colframe=codebg, rounded corners}
template <typename T>
void sort(T* array, size_t size);
\end{tcblisting}
\usemintedstyle{tango}
\item \alert{\textbf{Работает}} для примитивных типов, \alert{\textbf{есть}} проверки типов, \alert{\textbf{можно}} заинлайнить сравнение
\end{itemize}
\end{frame}

\begin{frame}[fragile]
\frametitle{Шаблоны}
\begin{itemize}
\usemintedstyle{lightbulb}
\begin{tcblisting}{listing engine=minted, minted language=cpp, minted options={fontsize=\scriptsize}, listing only, colback=codebg, colframe=codebg, rounded corners}
template <typename T>
void sort(T* array, size_t size);
\end{tcblisting}
\pause
\usemintedstyle{tango}
\item Для каждого типа \mintinline{cpp}|T| компилятор сгенерирует \textit{отдельную функцию}
\pause
\item \(\Longrightarrow\) В неё можно заинлайнить (подставить) код сравнения
\pause
\item \(\Longrightarrow\) Её можно оптимизировать отдельно, зная особенности типа \mintinline{cpp}|T|
\pause
\item \(\Longrightarrow\) Больше работы компилятору :)
\end{itemize}
\end{frame}

\begin{frame}[fragile]
\frametitle{Non-type template parameters}
\begin{itemize}
\item Аргументы шаблонов могут быть не только типами, но и значениями (примитивных типов; с C++20 -- простых структур)
\usemintedstyle{lightbulb}
\begin{tcblisting}{listing engine=minted, minted language=cpp, minted options={fontsize=\scriptsize}, listing only, colback=codebg, colframe=codebg, rounded corners}
template <size_t S>
class Bitset {
  uint8_t data[(S + 7) / 8];

public:
  bool get(size_t index) const {
    // Что-то с index/8 и index\%8
  }

  void set(size_t index, bool value) {
    // Что-то с index/8 и index\%8
  }
};
\end{tcblisting}
\end{itemize}
\end{frame}

\begin{frame}[fragile]
\frametitle{Non-type template parameters}
\begin{itemize}
\item Аргументы шаблонов могут быть не только типами, но и значениями (примитивных типов; с C++20 -- простых структур)
\usemintedstyle{lightbulb}
\begin{tcblisting}{listing engine=minted, minted language=cpp, minted options={fontsize=\scriptsize}, listing only, colback=codebg, colframe=codebg, rounded corners}
template <typename T, size_t N>
struct Vector {
  T coords[N];

  Vector operator* (T x) const {
    ...
  }

  Vector operator+ (const Vector<T, N>& v) const {
    ...
  }
};
\end{tcblisting}
\end{itemize}
\end{frame}

\begin{frame}[fragile]
\frametitle{Template template parameters}
\begin{itemize}
\item Хотим реализовать стек, но позволить пользователю выбрать структуру данных для хранения элементов:
\usemintedstyle{lightbulb}
\begin{tcblisting}{listing engine=minted, minted language=cpp, minted options={fontsize=\scriptsize}, listing only, colback=codebg, colframe=codebg, rounded corners}
template <typename T, typename Container>
struct Stack {
  Container data;

public:
  void push(const T& value);
  bool empty() const;
  T pop();
};

Stack<int, Array<int>> stack1;
Stack<double, List<double>> stack2;
\end{tcblisting}
\pause
\item В чём проблема? \pause Повторение кода + небезопасно:
\begin{tcblisting}{listing engine=minted, minted language=cpp, minted options={fontsize=\scriptsize}, listing only, colback=codebg, colframe=codebg, rounded corners}
// Потеря точности!
Stack<double, Array<int>> stack;
\end{tcblisting}
\end{itemize}
\end{frame}

\begin{frame}[fragile]
\frametitle{Template template parameters}
\begin{itemize}
\item Решение -- шаблонные шаблонные параметры:
\usemintedstyle{lightbulb}
\begin{tcblisting}{listing engine=minted, minted language=cpp, minted options={fontsize=\scriptsize}, listing only, colback=codebg, colframe=codebg, rounded corners}
template <typename T, template <typename> class Container>
struct Stack {
  Container<T> data;

public:
  ...
};

Stack<int, Array> stack1;
Stack<double, List> stack2;
\end{tcblisting}
\end{itemize}
\end{frame}

\begin{frame}[fragile]
\frametitle{Template template parameters}
\begin{itemize}
\item Можно сделать значение параметра по умолчанию:
\usemintedstyle{lightbulb}
\begin{tcblisting}{listing engine=minted, minted language=cpp, minted options={fontsize=\scriptsize}, listing only, colback=codebg, colframe=codebg, rounded corners}
template <typename T,
  template <typename> class Container = Deque
>
struct Stack { ... };

Stack<char*> stack; // Stack<char*, Deque>
\end{tcblisting}
\end{itemize}
\end{frame}

\begin{frame}[fragile]
\frametitle{Специализация шаблонов (полная)}
\begin{itemize}
\item Хотим сделать специфичное поведение для конкретного набора шаблонных параметров:
\usemintedstyle{lightbulb}
\begin{tcblisting}{listing engine=minted, minted language=cpp, minted options={fontsize=\scriptsize}, listing only, colback=codebg, colframe=codebg, rounded corners}
template <typename T>
class Array { ... };

template <>
class Array<bool> {
  // Оптимизация: храним bool'ы в битах,
  // а не отдельными объектами
  uint8_t* data;
  size_t size;

public:
  Array(size_t size)
    : data(new uint8_t[(size + 7) / 8])
    , size(size)
  {}

  ...
};
\end{tcblisting}
\end{itemize}
\end{frame}

\begin{frame}[fragile]
\frametitle{Специализация шаблонов (полная)}
\begin{itemize}
\item Хотим сделать специфичное поведение для конкретного набора шаблонных параметров:
\usemintedstyle{lightbulb}
\begin{tcblisting}{listing engine=minted, minted language=cpp, minted options={fontsize=\scriptsize}, listing only, colback=codebg, colframe=codebg, rounded corners}
template <typename T, size_t N>
struct Vector { ... };

template <>
struct Vector<float, 2> {
  // Оптимизация: используем SIMD intrinsics
  float coords __attribute__ ((vector_size(8)));
};

template <>
struct Vector<float, 4> {
  // Оптимизация: используем SIMD intrinsics
  float coords __attribute__ ((vector_size(16)));
};
\end{tcblisting}
\end{itemize}
\end{frame}

\begin{frame}[fragile]
\frametitle{Специализация шаблонов (частичная)}
\begin{itemize}
\item Частичная специализация -- для некоторого множетсва параметров:
\usemintedstyle{lightbulb}
\begin{tcblisting}{listing engine=minted, minted language=cpp, minted options={fontsize=\scriptsize}, listing only, colback=codebg, colframe=codebg, rounded corners}
template <typename T>
class Array<Array<T>> { ... };

template <size_t N>
struct Vector<float, N> { ... };
\end{tcblisting}
\end{itemize}
\end{frame}

\begin{frame}[fragile]
\frametitle{Обработка ошибок}
\begin{itemize}
\item Ошибки -- неотъемлемая часть программирования
\pause
\item Ошибки по вине программиста (передали не тот параметр, вызвали не тот метод, etc):
\pause
\begin{itemize}
\item Выявляют на стадии разработки / тестирования
\pause
\item В `идеальной' программе их не должно быть (но всегда есть)
\pause
\item В C и C++ их часто не обрабатывают
\end{itemize}
\pause
\item Ошибки по вине окружения / пользователя / входных данных (файл не существует, неверный формат ввода, etc)
\pause
\begin{itemize}
\item Могут произойти в любой программе
\pause
\item Их \textbf{обязательно} нужно обрабатывать!
\end{itemize}
\end{itemize}
\end{frame}

\begin{frame}[fragile]
\frametitle{Действия по обработке ошибки}
\begin{itemize}
\item Проверить наличие ошибки
\usemintedstyle{lightbulb}
\begin{tcblisting}{listing engine=minted, minted language=cpp, minted options={fontsize=\scriptsize}, listing only, colback=codebg, colframe=codebg, rounded corners}
if (!file) { ... }
\end{tcblisting}
\pause
\item Освободить ресурсы
\usemintedstyle{lightbulb}
\begin{tcblisting}{listing engine=minted, minted language=cpp, minted options={fontsize=\scriptsize}, listing only, colback=codebg, colframe=codebg, rounded corners}
delete[] buffer;
\end{tcblisting}
\pause
\item Сообщить пользователю / вызывающей функции
\usemintedstyle{lightbulb}
\begin{tcblisting}{listing engine=minted, minted language=cpp, minted options={fontsize=\scriptsize}, listing only, colback=codebg, colframe=codebg, rounded corners}
std::cerr << "Failed to open file\n";
\end{tcblisting}
\pause
\item Предпринять действия по восстановлению от ошибки
\usemintedstyle{lightbulb}
\begin{tcblisting}{listing engine=minted, minted language=cpp, minted options={fontsize=\scriptsize}, listing only, colback=codebg, colframe=codebg, rounded corners}
file = fopen("backup.data", "rb");
\end{tcblisting}
\end{itemize}
\end{frame}

\begin{frame}[fragile]
\frametitle{C-style обработка ошибок}
\begin{itemize}
\item Через возвращаемое значение:
\usemintedstyle{lightbulb}
\begin{tcblisting}{listing engine=minted, minted language=cpp, minted options={fontsize=\scriptsize}, listing only, colback=codebg, colframe=codebg, rounded corners}
FILE* fopen(const char* filename, const char* mode) {
  if (/* файл не найден */) {
    return NULL;
  }
  if (/* отказано в доступе */) {
    return NULL;
  }
  ...
}

int Array::get(int index) {
  if (index < 0) return -1;
}
\end{tcblisting}
\pause
\item Непонятно, что именно произошло и что с этим делать (`что-то пошло не так')
\pause
\item Не всегда есть разумное `ошибочное' значение возвращаемого типа
\end{itemize}
\end{frame}

\begin{frame}[fragile]
\frametitle{C-style обработка ошибок}
\begin{itemize}
\item Через возвращаемое значение + глобальный флаг:
\usemintedstyle{lightbulb}
\begin{tcblisting}{listing engine=minted, minted language=cpp, minted options={fontsize=\scriptsize}, listing only, colback=codebg, colframe=codebg, rounded corners}
FILE* fopen(const char* filename, const char* mode) {
  if (/* файл не найден */) {
    errno = ENOENT;
    return NULL;
  }
  if (/* отказано в доступе */) {
    errno = EPERM;
    return NULL;
  }
  ...
}
\end{tcblisting}
\pause
\usemintedstyle{tango}
\item Нужно следить, чтобы другая функция не перезаписала \mintinline{cpp}|errno|
\pause
\item Нельзя передать более подробную информацию об ошибке
\end{itemize}
\end{frame}

\begin{frame}[fragile]
\frametitle{C-style обработка ошибок}
\begin{itemize}
\item Через возвращаемое значение + дополнительный аргумент:
\usemintedstyle{lightbulb}
\begin{tcblisting}{listing engine=minted, minted language=cpp, minted options={fontsize=\scriptsize}, listing only, colback=codebg, colframe=codebg, rounded corners}
FILE* fopen(const char* filename, const char* mode, int* error)
{
  if (/* файл не найден */) {
    *error = ENOENT;
    return NULL;
  }
  if (/* отказано в доступе */) {
    *error = EPERM;
    return NULL;
  }
  ...
}
\end{tcblisting}
\pause
\usemintedstyle{tango}
\item Перемешивается основная логика функции и обработка ошибок
\end{itemize}
\end{frame}

\begin{frame}[fragile]
\frametitle{Result / Either / std::expected}
\begin{itemize}
\item Современный подход -- использовать тип-сумму:
\usemintedstyle{lightbulb}
\begin{tcblisting}{listing engine=minted, minted language=cpp, minted options={fontsize=\scriptsize}, listing only, colback=codebg, colframe=codebg, rounded corners}
std::expected<FILE*, int> fopen(
  const char* filename, const char* mode)
{
  if (/* файл не найден */) {
    return std::unexpected(ENOENT);
  }
  if (/* отказано в доступе */) {
    return std::unexpected(EPERM);
  }
  FILE* file = ...;
  ...
  return file;
}
\end{tcblisting}
\pause
\usemintedstyle{tango}
\item Самый аккуратный подход
\pause
\item Всё ещё перемешивается основная логика функции и обработка ошибок
\end{itemize}
\end{frame}

\begin{frame}[fragile]
\frametitle{Исключения}
\begin{itemize}
\usemintedstyle{lightbulb}
\begin{tcblisting}{listing engine=minted, minted language=cpp, minted options={fontsize=\scriptsize}, listing only, colback=codebg, colframe=codebg, rounded corners}
class DivideByZeroException {
public:
  const char* what() {
    return "division by zero";
  }
};

double divide(int a, int b) {
  if (b == 0) {
    throw DivideByZeroException();
  }
  return double(a) / b;
}
\end{tcblisting}
\end{itemize}
\end{frame}

\begin{frame}[fragile]
\frametitle{Исключения}
\begin{itemize}
\usemintedstyle{lightbulb}
\begin{tcblisting}{listing engine=minted, minted language=cpp, minted options={fontsize=\scriptsize}, listing only, colback=codebg, colframe=codebg, rounded corners}
try {
  c = divide(a, b);
}
catch (DivideByZeroException& e) {
  // Решаем, что с этим делать
}
\end{tcblisting}
\pause
\usemintedstyle{tango}
\item \mintinline{cpp}|throw| \textit{бросает} исключение
\pause
\item \mintinline{cpp}|catch| \textit{ловит} исключение
\end{itemize}
\end{frame}

\begin{frame}[fragile]
\frametitle{Исключения}
\begin{itemize}
\usemintedstyle{tango}
\item Брошенное исключение идёт вверх по стеку вызовов, пока не наткнётся на подходящий \mintinline{cpp}|try/catch| блок \textit{(раскрутка стека -- stack unwinding)}:
\usemintedstyle{lightbulb}
\begin{tcblisting}{listing engine=minted, minted language=cpp, minted options={fontsize=\scriptsize}, listing only, colback=codebg, colframe=codebg, rounded corners}
int f(int x) {
  return divide(1000, x);
}

int g(int x) {
  return f(x - 42);
}

int main() {
  try {
    g(42);
  }
  catch (const DivideByZeroException& e) {
    ...
  }
}
\end{tcblisting}
\end{itemize}
\end{frame}

\begin{frame}[fragile]
\frametitle{Исключения}
\begin{itemize}
\usemintedstyle{tango}
\item После бросания исключения нормальный поток исполнения останавливается -- \mintinline{cpp}|printf| не будет вызван:
\usemintedstyle{lightbulb}
\begin{tcblisting}{listing engine=minted, minted language=cpp, minted options={fontsize=\scriptsize}, listing only, colback=codebg, colframe=codebg, rounded corners}
void f() {
  throw Exception(...);
  printf("I am Groot");
}
\end{tcblisting}
\pause
\usemintedstyle{tango}
\item Если подходящего \mintinline{cpp}|try/catch| блока не нашлось, программа аварийно завершается:
\usemintedstyle{lightbulb}
\begin{tcblisting}{listing engine=minted, minted language=cpp, minted options={fontsize=\scriptsize}, listing only, colback=codebg, colframe=codebg, rounded corners}
int main() {
  // Нет try/catch
  throw Exception(...);
}
\end{tcblisting}
\end{itemize}
\end{frame}

\begin{frame}[fragile]
\frametitle{Несколько типов исключений}
\begin{itemize}
\usemintedstyle{tango}
\item Если вызываемый код может бросать несколько типов исключений, можно обработать их все:
\usemintedstyle{lightbulb}
\begin{tcblisting}{listing engine=minted, minted language=cpp, minted options={fontsize=\scriptsize}, listing only, colback=codebg, colframe=codebg, rounded corners}
int main() {
  try {
    runServer();
  }
  catch (const FilesystemError& e) {
    ...
  }
  catch (const NetworkError& e) {
    ...
  }
  catch (const InternalError& e) {
    ...
  }
}
\end{tcblisting}
\end{itemize}
\end{frame}

\begin{frame}[fragile]
\frametitle{Иерархия исключений}
\begin{itemize}
\usemintedstyle{tango}
\item \mintinline{cpp}|catch (Type)| поймает не только тип \mintinline{cpp}|Type|, но и любого его наследника:
\usemintedstyle{lightbulb}
\begin{tcblisting}{listing engine=minted, minted language=cpp, minted options={fontsize=\scriptsize}, listing only, colback=codebg, colframe=codebg, rounded corners}
class Exception {
public:
  virtual const char* what() { return "Error"; }
};

class DivideByZeroException : public Exception { ... };

int main() {
  try {
    divide(10, 0);
  }
  catch (const Exception& e) {
    std::cerr << e.what() << std::endl;
  }
}
\end{tcblisting}
\end{itemize}
\end{frame}

\begin{frame}[fragile]
\frametitle{Иерархия исключений}
\begin{itemize}
\usemintedstyle{tango}
\item Часто делают иерархии классов исключений для более точной обработки ошибок:
\usemintedstyle{lightbulb}
\begin{tcblisting}{listing engine=minted, minted language=cpp, minted options={fontsize=\scriptsize}, listing only, colback=codebg, colframe=codebg, rounded corners}
class Exception { ... };

class FilesystemError   : public Exception { ... };
class FileNotFoundError : public FilesystemError { ... };
class PermissionDenied  : public FilesystemError { ... };

class NetworkError      : public Exception { ... };
class ConnectionLost    : public NetworkError { ... };
class ConnectionTimeout : public NetworkError { ... };

class InternalError     : public Exception { ... };
class AssertionFailed   : public InternalError { ... };
class UnreachableCode   : public InternalError { ... };
\end{tcblisting}
\pause
\usemintedstyle{tango}
\item В стандартной библиотеке все исключения -- наследники \mintinline{cpp}|std::exception|
\end{itemize}
\end{frame}

\begin{frame}[fragile]
\frametitle{Иерархия исключений}
\begin{itemize}
\usemintedstyle{tango}
\item Важно: в списке \mintinline{cpp}|catch| блоков выбирается \textbf{первый подходящий}!
\usemintedstyle{lightbulb}
\begin{tcblisting}{listing engine=minted, minted language=cpp, minted options={fontsize=\scriptsize}, listing only, colback=codebg, colframe=codebg, rounded corners}
try {
  ...
}
catch (const Exception& e) {
  ...
}
catch (const ConnectionLost& e) {
  // Никогда не вызовется!
}
\end{tcblisting}
\end{itemize}
\end{frame}

\begin{frame}[fragile]
\frametitle{Иерархия исключений}
\begin{itemize}
\usemintedstyle{tango}
\item Правильный порядок:
\usemintedstyle{lightbulb}
\begin{tcblisting}{listing engine=minted, minted language=cpp, minted options={fontsize=\scriptsize}, listing only, colback=codebg, colframe=codebg, rounded corners}
try {
  ...
}
catch (const ConnectionLost& e) {
  // Сначала наследник,
}
catch (const Exception& e) {
  // потом базовый класс исключения
}
\end{tcblisting}
\end{itemize}
\end{frame}

\begin{frame}[fragile]
\frametitle{Исключения}
\begin{itemize}
\usemintedstyle{tango}
\item Можно поймать исключение неизвестного типа конструкией \mintinline{cpp}|catch(...)| (но тогда с ним ничего особо не сделать):
\usemintedstyle{lightbulb}
\begin{tcblisting}{listing engine=minted, minted language=cpp, minted options={fontsize=\scriptsize}, listing only, colback=codebg, colframe=codebg, rounded corners}
try {
  runServer();
}
catch (...) {
  std::cerr << "Unknown error" << std::endl;
}
\end{tcblisting}
\pause
\item \textbf{NB:} Исключение может быть любым типом, но лучше так не делать :)
\usemintedstyle{lightbulb}
\begin{tcblisting}{listing engine=minted, minted language=cpp, minted options={fontsize=\scriptsize}, listing only, colback=codebg, colframe=codebg, rounded corners}
throw NULL;
throw 42;
throw "Oopsie woopsie";
\end{tcblisting}
\end{itemize}
\end{frame}

\begin{frame}[fragile]
\frametitle{Исключения}
\begin{itemize}
\usemintedstyle{tango}
\item Можно бросить пойманное исключение дальше \textit{(rethrow)} после его обработки конструкией \mintinline{cpp}|throw;| (даже если оно неизвестного типа):
\usemintedstyle{lightbulb}
\begin{tcblisting}{listing engine=minted, minted language=cpp, minted options={fontsize=\scriptsize}, listing only, colback=codebg, colframe=codebg, rounded corners}
try {
  runServer();
}
catch (const NetworkError& e) {
  std::cerr << "Network error: " << e.what() << std::endl;
  throw;
}
catch (...) {
  std::cerr << "Unknown error" << std::endl;
  throw;
}
\end{tcblisting}
\end{itemize}
\end{frame}

\begin{frame}[fragile]
\frametitle{Исключения и деструкторы}
\begin{itemize}
\usemintedstyle{tango}
\item В процессе раскрутки стека автоматически вызываются деструкторы всех локальных переменных:
\usemintedstyle{lightbulb}
\begin{tcblisting}{listing engine=minted, minted language=cpp, minted options={fontsize=\scriptsize}, listing only, colback=codebg, colframe=codebg, rounded corners}
void f() {
  Array array(100);
  if (...) {
    // Вызовется array.~Array()
    throw Exception(...);
  }
}

void g() {
  File file(...);
  // Если бросится исключение, вызовется file.~File()
  f();
}
\end{tcblisting}
\end{itemize}
\end{frame}

\begin{frame}[fragile]
\frametitle{Исключения и деструкторы}
\begin{itemize}
\usemintedstyle{lightbulb}
\begin{tcblisting}{listing engine=minted, minted language=cpp, minted options={fontsize=\scriptsize}, listing only, colback=codebg, colframe=codebg, rounded corners}
void f() {
  int* array = new int[100];
  if (...) {
    // Формально, вызовется деструктор array
    // Это примитивный тип, так что деструктор ничего не делает
    // => Утечка памяти
    throw Exception(...);
  }
  delete[] array;
}
\end{tcblisting}
\end{itemize}
\end{frame}

\begin{frame}[fragile]
\frametitle{Идиома RAII}
\begin{itemize}
\usemintedstyle{tango}
\item \textbf{Resource Acquisition Is Initialization} -- взятие ресурса должно происходить через инициализацию
\pause
\item По существу RAII обычно сводится к освобождению ресурса в деструкторе
\pause
\item Не сможем забыть очистить ресурсы\pause, теперь и в случае исключений!
\usemintedstyle{lightbulb}
\begin{tcblisting}{listing engine=minted, minted language=cpp, minted options={fontsize=\scriptsize}, listing only, colback=codebg, colframe=codebg, rounded corners}
void f() {
  Array array(100);
  if (...) {
    // Вызовется array.~Array()
    // Утечки памяти нет
    throw Exception(...);
  }
  delete[] array;
}
\end{tcblisting}
\end{itemize}
\end{frame}

\begin{frame}[fragile]
\frametitle{Исключения в конструкторе}
\begin{itemize}
\usemintedstyle{tango}
\item До конца работы конструктора объекта, объект не считается проинициализированным, и его деструктор не вызовется:
\usemintedstyle{lightbulb}
\begin{tcblisting}{listing engine=minted, minted language=cpp, minted options={fontsize=\scriptsize}, listing only, colback=codebg, colframe=codebg, rounded corners}
// Может бросить исключение
void checkSupported(const char *effect);

class AudioEffect {
  float* buffer;

public:
  AudioEffect(const char* effect) {
    buffer = new float[1024];
    // Деструктор ~AudioEffect не вызовется
    // => Утечка памяти
    checkSupported(effect);
  }

  ~AudioEffect() {
    delete[] buffer;
  }
};
\end{tcblisting}
\end{itemize}
\end{frame}

\begin{frame}[fragile]
\frametitle{Исключения в конструкторе}
\begin{itemize}
\usemintedstyle{tango}
\item Решение №1:
\usemintedstyle{lightbulb}
\begin{tcblisting}{listing engine=minted, minted language=cpp, minted options={fontsize=\scriptsize}, listing only, colback=codebg, colframe=codebg, rounded corners}
  AudioEffect(const char* effect) {
    buffer = new float[1024];
    try {
      checkSupported(effect);
    }
    catch (...) {
      delete[] buffer;
      throw;
    }
  }
\end{tcblisting}
\end{itemize}
\end{frame}

\begin{frame}[fragile]
\frametitle{Исключения в конструкторе}
\begin{itemize}
\item Решение №2 -- RAII:
\usemintedstyle{lightbulb}
\begin{tcblisting}{listing engine=minted, minted language=cpp, minted options={fontsize=\scriptsize}, listing only, colback=codebg, colframe=codebg, rounded corners}
class AudioEffect {
  Array<float> buffer;

public:
  AudioEffect(const char* effect)
    : buffer(1024)
  {
    // this->buffer уже проинициализирован
    // Деструктор buffer.~Array() будет вызван
    checkSupported(effect);
  }
};
\end{tcblisting}
\end{itemize}
\end{frame}

\begin{frame}[fragile]
\frametitle{Исключения в конструкторе}
\begin{itemize}
\item Что, если конструктор одного из полей бросает исключение?
\usemintedstyle{lightbulb}
\begin{tcblisting}{listing engine=minted, minted language=cpp, minted options={fontsize=\scriptsize}, listing only, colback=codebg, colframe=codebg, rounded corners}
class AudioManager {
  FILE* input;
  AudioEffect effect;

public:
  // Если конструктор effect бросит исключение,
  // деструктор AudioManager не вызовется
  // и файл input не закроется
  AudioManager()
    : input(fopen("meow.wav", "rb"))
    , effect("wrum-wrum")
  {}

  ~AudioManager() {
    fclose(input);
  }
};
\end{tcblisting}
\end{itemize}
\end{frame}

\begin{frame}[fragile]
\frametitle{Исключения в конструкторе}
\begin{itemize}
\item Решение №1 -- особый синтаксис \mintinline{cpp}|try| блока:
\usemintedstyle{lightbulb}
\begin{tcblisting}{listing engine=minted, minted language=cpp, minted options={fontsize=\scriptsize}, listing only, colback=codebg, colframe=codebg, rounded corners}
AudioManager() try
  : input(fopen("meow.wav", "rb"))
  , effect("wrum-wrum")
{
  // тело конструктора
}
catch (...) {
  fclose(input);
  throw;
}
\end{tcblisting}
\pause
\item \textbf{NB}: работает и для функций
\usemintedstyle{lightbulb}
\begin{tcblisting}{listing engine=minted, minted language=cpp, minted options={fontsize=\scriptsize}, listing only, colback=codebg, colframe=codebg, rounded corners}
int main() try {
  // тело функции
}
catch (const Exception& e) {
  std::cerr << "Error: " << e.what() << std::endl;
}
\end{tcblisting}
\end{itemize}
\end{frame}

\begin{frame}[fragile]
\frametitle{Исключения в конструкторе}
\begin{itemize}
\item Решение №2 -- RAII:
\usemintedstyle{lightbulb}
\begin{tcblisting}{listing engine=minted, minted language=cpp, minted options={fontsize=\scriptsize}, listing only, colback=codebg, colframe=codebg, rounded corners}
class AudioManager {
  FileStream input;
  AudioEffect effect;

public:
  // Если конструктор effect бросит исключение,
  // вызовется деструктор input
  AudioManager()
    : input("meow.wav")
    , effect("wrum-wrum")
  {}
};
\end{tcblisting}
\end{itemize}
\end{frame}

\begin{frame}[fragile]
\frametitle{RAII}
\begin{itemize}
\item Совет: если вам приходится вручную обрабатывать исключения в конструкторе, скорее всего, какие-то поля надо выделить в отдельные RAII-классы :)
\end{itemize}
\end{frame}

\begin{frame}[fragile]
\frametitle{Исключения в деструкторе}
\begin{itemize}
\usemintedstyle{lightbulb}
\begin{tcblisting}{listing engine=minted, minted language=cpp, minted options={fontsize=\scriptsize}, listing only, colback=codebg, colframe=codebg, rounded corners}
class Database {
  ...
  ~Database() {
    // Может бросить LoggingError
    logger.log("Database is closed");
  }
};
\end{tcblisting}
\pause
\usemintedstyle{tango}
\item Деструктор \mintinline{cpp}|~Database| может вызваться в процессе раскрутки стека из-за другого исключения
\pause
\item В других языках часто поддерживаются цепочки исключений (\mintinline{cpp}|LoggingError caused by DatabaseError|)
\pause
\item В C++ исключения могут быть произвольного типа \(\Longrightarrow\) в общем случае к ним не добавить никакую информацию
\pause
\item Поэтому в C++ исключение, брошенное в деструкторе, приводит к аварийному завершению программы
\end{itemize}
\end{frame}

\begin{frame}[fragile]
\frametitle{Исключения в деструкторе}
\begin{itemize}
\item Решение №1 -- проигнорировать исключение:
\usemintedstyle{lightbulb}
\begin{tcblisting}{listing engine=minted, minted language=cpp, minted options={fontsize=\scriptsize}, listing only, colback=codebg, colframe=codebg, rounded corners}
class Database {
  ...
  ~Database() {
    try {
      logger.log("Database is closed");
    }
    catch (...) {}
  }
};
\end{tcblisting}
\pause
\item Плохо, потому что мы не узнаем о какой-то ошибке (но так всё равно часто делают)
\end{itemize}
\end{frame}

\begin{frame}[fragile]
\frametitle{Исключения в деструкторе}
\begin{itemize}
\item Решение №2 -- вынести такой код из деструктора и вызывать явно:
\usemintedstyle{lightbulb}
\begin{tcblisting}{listing engine=minted, minted language=cpp, minted options={fontsize=\scriptsize}, listing only, colback=codebg, colframe=codebg, rounded corners}
class Database {
  ...
  void close() {
    logger.log("Database is closed");
  }

  ~Database() {}
};

int main() {
  Database db;
  ...
  db.close();
}
\end{tcblisting}
\pause
\usemintedstyle{tango}
\item Плохо, потому что можно забыть вызвать \mintinline{cpp}|close|
\end{itemize}
\end{frame}

\begin{frame}[fragile]
\frametitle{Исключения в деструкторе}
\begin{itemize}
\item Есть ли хорошее решение? \pause Нет :(
\end{itemize}
\end{frame}

\begin{frame}[fragile]
\frametitle{Гарантии при работе с исключениями}
\begin{itemize}
\item Часто функции/методы в коде на C++ предоставляют некоторые гарантии в связи с исключениями (обычно это описано в документации)
\pause
\item \textbf{Nothrow guarantee}: функция/метод не бросает исключение
\pause
\item \textbf{Strong guarantee}: в случае исключения состояние программы откатывается до состояния перед вызовом функции/метода
\pause
\item \textbf{Basic guarantee}: в случае исключения программа останется в некотором допустимом состоянии
\pause
\item \textbf{No guarantee}: в случае исключения \shrug
\end{itemize}
\end{frame}

\begin{frame}[fragile]
\frametitle{No guarantee}
\begin{itemize}
\usemintedstyle{lightbulb}
\begin{tcblisting}{listing engine=minted, minted language=cpp, minted options={fontsize=\scriptsize}, listing only, colback=codebg, colframe=codebg, rounded corners}
void process() {
  int* buffer = new int[1024];
  FILE* file = fopen(...);
  if (!file) {
    // Утечка памяти
    throw Exception(...);
  }
}
\end{tcblisting}
\pause
\item Самый плохой вариант: обычно стараются давать хоть какую-нибудь гарантию, иначе кодом невозможно пользоваться корректно
\end{itemize}
\end{frame}

\begin{frame}[fragile]
\frametitle{Basic guarantee}
\begin{itemize}
\usemintedstyle{lightbulb}
\begin{tcblisting}{listing engine=minted, minted language=cpp, minted options={fontsize=\scriptsize}, listing only, colback=codebg, colframe=codebg, rounded corners}
void update(Array<int>& array) {
  for (size_t i = 0; i < array.size(); ++i) {
    // Может бросить FileNotFoundError
    FileStream file(getFilename(i));
    array[i] = file.readNextInt();
  }
}
\end{tcblisting}
\pause
\item Слабая гарантия: состояние программы останется валидным (нет утечек памяти, битых ссылок, и т.п.), но уже затронутые элементы массива останутся изменёнными
\pause
\item Обычно легко достигается грамотным применением RAII
\end{itemize}
\end{frame}

\begin{frame}[fragile]
\frametitle{Strong guarantee}
\begin{itemize}
\usemintedstyle{lightbulb}
\begin{tcblisting}{listing engine=minted, minted language=cpp, minted options={fontsize=\scriptsize}, listing only, colback=codebg, colframe=codebg, rounded corners}
void update(Array<int>& array) {
  Array<int> newArray(array.size());

  for (size_t i = 0; i < array.size(); ++i) {
    // Может бросить FileNotFoundError
    FileStream file(getFilename(i));
    newArray[i] = file.readNextInt();
  }
  
  // Идиома copy-and-swap
  array.swap(newArray);
}
\end{tcblisting}
\pause
\item Сильная гарантия: состояние программы не изменится
\pause
\item Обычно тяжело обеспечить
\pause
\item Стандартные контейнеры стараются дать сильную гарантию везде, где это возможно
\end{itemize}
\end{frame}

\begin{frame}[fragile]
\frametitle{Nothrow guarantee}
\begin{itemize}
\usemintedstyle{lightbulb}
\begin{tcblisting}{listing engine=minted, minted language=cpp, minted options={fontsize=\scriptsize}, listing only, colback=codebg, colframe=codebg, rounded corners}
int strlen(const char* str) {
  int result = 0;
  for (; *s != 0; ++str, ++result);
  return result;
}

vec2f rotate(vec2f vec, float angle) {
  float c = cos(angle);
  float s = sin(angle);
  return vec2f(
    c * vec.x - s * vec.y,
    s * vec.x + c * vec.y,
  );
}
\end{tcblisting}
\pause
\item Самая сильная гарантия: обычно означает, что код в принципе никак не может бросить исплючение
\end{itemize}
\end{frame}

\end{document}